\documentclass[12pt]{article}

% ---------- packages ----------
\usepackage[margin=1in]{geometry}
\usepackage{setspace}
\usepackage{amsmath}
\usepackage{amssymb}
\usepackage{graphicx}
\usepackage{booktabs}
\usepackage{caption}
\usepackage{subcaption}
\usepackage{float}
\usepackage{hyperref}
\usepackage{natbib}
\usepackage{microtype}
\usepackage{lmodern}
\usepackage[T1]{fontenc}
\usepackage[utf8]{inputenc}
\usepackage{array}
\usepackage{tabularx}
\usepackage{longtable}

% ---------- hyperref setup ----------
\hypersetup{
  colorlinks=true,
  linkcolor=black,
  citecolor=black,
  urlcolor=black
}

% ---------- spacing ----------
\doublespacing

% ---------- title block ----------
\title{\textbf{Wage Transparency and Within-Rank Salary Compression:\\
Evidence from UC Faculty Compensation, 2012--2024}}

\author{
  Joshua Kenworthy \and
  Agnibha Bhattacharya \and
  Jackson Wurzer
}

\date{March 2026}

% ============================================================
\begin{document}
\maketitle
\thispagestyle{empty}

% ---------- Abstract ----------
\begin{abstract}
\noindent
We examine whether the 2022 increase in the salience of wage information
in the University of California (UC) system corresponds with measurable
changes in within-rank faculty salary dispersion. Using a 13-year
repeated cross-section of approximately 229,000 salary records from
2012 to 2024---covering five faculty ranks across ten UC campuses---we
construct annual coefficients of variation (CV) for each rank and
estimate log-linear OLS regressions with rank fixed effects, campus
fixed effects, and a linear time trend. We find that mean nominal
salaries grew by roughly 52\,\% over the sample period before declining
6.4\,\% in 2024. The baseline regression yields a \textit{post2022}
coefficient of $-0.063$ ($p < 0.001$), indicating wages were
approximately 6.3\,\% below the secular trend after 2022 conditional on
rank and campus. The CV declined modestly for Full, Associate, Research,
and Clinical professors after 2022, with Associate professors showing
the largest drop of approximately 4.6 percentage points. Assistant
professors are an exception: their CV rose by roughly 1.4 percentage
points, consistent with external labor-market pressures. Levene's tests
confirm that absolute variance increased significantly across all ranks
post-2022, but this is reconcilable with CV compression because mean
salaries rose faster than the absolute spread. A
difference-in-differences (DiD) interaction model finds no significant
narrowing of the wage gap between Full and Assistant professors, though
Associate and Clinical professors fared relatively better in the
post-period. We interpret these patterns as consistent with
transparency-driven internal equity adjustments concentrated within,
rather than across, faculty ranks.
\end{abstract}

\newpage
\tableofcontents
\newpage

% ============================================================
\section{Introduction}
\label{sec:intro}

The University of California (UC) system is one of the largest public
employers in the United States, employing over 265,000 faculty and staff
across ten campuses \citep{UC2025}. Compensation within the UC system has
long been subject to a high degree of transparency. Following a 2008
court decision, the \textit{Sacramento Bee} created a searchable online
database that allowed the public to view the salary of any California
state employee, including UC faculty \citep[p.~2982]{card2012}. Since
that time, multiple websites have made UC salary data easily accessible,
enabling employees to compare their compensation with that of their peers
on an ongoing basis.

This institutional setting provides a useful environment to study the
relationship between wage transparency and wage compression. Economic
theory predicts that when workers can observe peer compensation,
within-group pay disparities become salient and can generate
dissatisfaction, reduce effort, or increase turnover among lower-paid
employees \citep[p.~2981]{card2012}. Employers may therefore face
incentives to narrow wage gaps within comparable job categories in order
to retain workers and maintain productivity \citep{cullen2022}.

Our paper focuses on the period surrounding 2022, which represents a
meaningful shift in the compensation and transparency environment for UC
faculty. Two developments converged: first, California's Pay
Transparency Act took effect in January 2023, expanding wage disclosure
requirements across the state; second, in late 2022 the UC system
experienced its largest academic labor action in history, with tens of
thousands of academic workers striking over wages and working conditions
and obtaining significant contract settlements. Together, these events
substantially elevated the salience of wage comparisons within the UC
system---even though salary data had technically been public for years.

We examine whether this period corresponds with measurable changes in
within-rank salary dispersion among UC faculty. Our analysis focuses on
five faculty categories: Assistant, Associate, Full, Clinical, and
Research professors. The main empirical tools are: (1) annual
coefficients of variation (CV) computed separately for each rank; (2) a
baseline log-linear OLS regression that estimates whether log wages
deviated from the secular trend after 2022; (3) Levene's tests for
equality of variance across the pre- and post-2022 periods; and (4) a
difference-in-differences interaction model that tests whether wage
gaps between ranks changed after 2022.

The remainder of the paper proceeds as follows. Section~\ref{sec:theory}
describes the theoretical framework. Section~\ref{sec:data} describes
the data sources and preprocessing. Section~\ref{sec:methods} outlines
the empirical methodology. Section~\ref{sec:results} presents the
results. Section~\ref{sec:discussion} discusses findings and limitations.
Section~\ref{sec:conclusion} concludes. Figures and supplementary tables
appear in the Appendix.

% ============================================================
\section{Theoretical Framework}
\label{sec:theory}

We model the 2022 transparency and labor-action shock as an event that
abruptly increased the salience of peer wage comparisons within the UC
faculty labor market. The theoretical mechanism draws on two strands of
the economics literature.

The first strand concerns social preferences and inequality aversion.
\citet{bandiera2005} and \citet{card2012} document that workers who
learn they are paid less than comparable peers experience a significant
reduction in job satisfaction and effort. Following the utility
specification in \citet{bandiera2005}, an agent's utility can be
written as:
\begin{equation}
  U_i = W_i - C(e_i) - \lambda\max(0,\; W_{\text{peer}} - W_i),
  \label{eq:utility}
\end{equation}
where $W_i$ is the worker's wage, $C(e_i)$ is the cost of effort, and
the final term captures the disutility---an ``envy cost''---from
observing a higher-paid peer. The parameter $\lambda \geq 0$ indexes the
salience of wage comparisons. When compensation is opaque, $\lambda
\approx 0$. The 2022 transparency shock raises $\lambda$, shifting the
agent's incentive-compatibility constraints and creating pressure for the
principal (the UC system) to reduce within-group wage dispersion in order
to prevent a decline in aggregate effort.

The second strand concerns the equilibrium effects of wage transparency.
\citet{cullen2022} show that transparency fundamentally alters bargaining
dynamics by making individualized pay differences highly visible, and
consequently costly for firms to sustain. In equilibrium, employers
facing transparent wage structures rationally choose to reduce
within-group dispersion to limit internal dissatisfaction. An important
boundary condition in the Cullen--Pakzad-Hurson framework is that this
compression is primarily an internal equity response: it operates within
peer groups defined by comparable job classification, not across
structurally different positions.

Applied to the UC setting, this framework generates two predictions. For
established faculty ranks (Full, Associate, Clinical, Research), where
salaries are largely set through internal processes and promotion
ladders, we expect the transparency shock to induce compression: the CV
should fall and, potentially, rank-level wages should lie below their
pre-trend trajectory. For entry-level Assistant professors, however,
salaries are primarily determined by competitive external hiring, so the
UC system has limited ability to compress wages without risking
recruitment failure. We therefore predict little or no compression at
the bottom of the hierarchy.

% ============================================================
\section{Data}
\label{sec:data}

\subsection{Sources}

Our salary data come from two public administrative databases. Data for
2012 are drawn from the UC Annual Wage Database maintained by the UC
Office of the President (\url{https://ucannualwage.ucop.edu/wage/}).
Data for 2013 through 2024 are drawn from the California State
Controller's Government Compensation in California database
(\url{https://publicpay.ca.gov/}). All data are publicly accessible and
contain no individual identifiers.

The merged panel covers 13 annual snapshots (2012--2024) and, after
preprocessing, contains 229,386 professor-year observations with
positive total wages, and 225,412 observations with positive regular pay.

\subsection{Preprocessing and Sample Construction}

Each annual file was loaded and standardized to five core columns: year,
employer name, position title, regular pay, and total wages. We retained
only records whose position title matched a ``professor'' pattern using
a case-insensitive regular expression that captures titles containing
``Prof'' followed by a space or hyphen, thereby including all professor
variants while excluding ``Professional'' administrative staff.

We then mapped free-text position titles into five faculty ranks using a
priority-ordered keyword rule: a position is classified as
\textit{Assistant} if the title contains ``assistant'' or ``asst'';
\textit{Associate} if it contains ``associate'' or ``assoc'';
\textit{Clinical} if it contains ``clinical'' or ``clin'';
\textit{Research} if it contains ``research'' or ``res''; and
\textit{Full} otherwise (including bare ``Professor'' titles). This
within-rank classification is central to our identification strategy:
comparing an entry-level Assistant Professor to a tenured Full Professor
would confound structural position differences with salary inequality. By
analyzing dispersion separately within each rank, we isolate
within-peer-group variance.

Before finalizing the panel, we removed duplicate rows, dropped records
with negative or zero pay values, and verified that all five required
columns were non-missing. The resulting analysis dataset is a
\textit{repeated cross-section}: annual CSVs contain no individual
identifiers, so we cannot track the same professor across years. Each
observation is an independent salary record for whoever held a given
position in a given year.

\subsection{Descriptive Statistics}

Table~\ref{tab:trends} summarizes the annual sample size and mean salary.
The total number of professor observations grew from 16,017 in 2012 to a
peak of 20,176 in 2023 before declining to 18,898 in 2024.
Figure~\ref{fig:salary_trends} plots the level and distributional
evolution of nominal total wages.

\begin{table}[H]
\centering
\caption{Year-by-Year Professor Observations and Mean Total Wages, 2012--2024}
\label{tab:trends}
\begin{tabular}{lrrrrr}
\toprule
Year & N & Mean (\$) & Median (\$) & Std.\ Dev.\ (\$) & YoY\,\% \\
\midrule
2012 & 16,017 & 162,636 & 138,400 & 117,478 & --- \\
2013 & 15,711 & 172,838 & 146,505 & 122,877 & +6.3\% \\
2014 & 16,042 & 179,202 & 152,000 & 128,030 & +3.7\% \\
2015 & 16,508 & 187,178 & 158,688 & 133,957 & +4.5\% \\
2016 & 17,176 & 195,355 & 165,121 & 142,233 & +4.4\% \\
2017 & 17,739 & 202,788 & 171,252 & 148,112 & +3.8\% \\
2018 & 18,151 & 210,399 & 178,817 & 151,300 & +3.8\% \\
2019 & 17,963 & 218,634 & 185,050 & 157,137 & +3.9\% \\
2020 & 19,224 & 225,799 & 195,000 & 154,026 & +3.3\% \\
2021 & 19,415 & 237,191 & 201,700 & 166,500 & +5.0\% \\
2022 & 19,737 & 249,471 & 211,045 & 174,503 & +5.2\% \\
2023 & 20,176 & 264,251 & 224,633 & 186,085 & +5.9\% \\
2024 & 18,898 & 247,276 & 208,691 & 179,869 & $-6.4$\% \\
\bottomrule
\end{tabular}
\begin{flushleft}
\small Source: UC Annual Wage Database (2012) and California State
Controller's Office (2013--2024). Sample restricted to professor
positions with positive total wages.
\end{flushleft}
\end{table}

Mean nominal wages grew steadily from \$162,636 in 2012 to a peak of
\$264,251 in 2023---a cumulative increase of approximately 52\,\%, or
roughly 4\,\% per year on average. Year-over-year growth was remarkably
stable from 2014 through 2021, ranging from 3.3\,\% to 5.0\,\%
annually, suggesting a consistent underlying wage-setting regime during
the pre-period. The acceleration in 2022 (+5.2\,\%) and 2023 (+5.9\,\%)
stands out modestly above that trend, as does the sharp reversal in 2024
($-6.4$\,\%).

Across ranks in 2024, Clinical faculty hold the highest mean compensation
at approximately \$371,264 for the rank overall, reflecting the medical
compensation structure at UC health science campuses. Full Professors
earn a mean of \$261,329, Associate Professors \$233,613, and Assistant
Professors \$183,209. Among campuses, UC San Francisco (\$298,740) and
UC San Diego (\$292,983) report the highest mean total wages, while UC
Riverside (\$167,993) and UC Merced (\$164,933) report the lowest.

All analysis uses nominal wages. No inflation adjustment is applied;
the analysis focuses on within-period distributional patterns rather
than real purchasing power comparisons.

% ============================================================
\section{Empirical Strategy}
\label{sec:methods}

\subsection{Measuring Within-Rank Wage Dispersion}

Our primary measure of wage compression is the coefficient of variation
(CV), defined as the ratio of the standard deviation to the mean within
a rank-year cell:
\begin{equation}
  \text{CV}_{r,t} = \frac{\sigma_{r,t}}{\mu_{r,t}} \times 100,
\end{equation}
where $r$ indexes rank and $t$ indexes year. The CV normalizes
dispersion by the mean level, allowing meaningful comparison across
periods with different average salary levels. Because nominal wages rose
substantially over the sample period, a raw comparison of standard
deviations would be uninformative; the CV controls for this scale effect.
We also report the standard deviation and interquartile range (IQR) as
supplementary measures.

\subsection{Baseline OLS Regression}

To estimate the conditional association between the post-2022 period and
log wages, we estimate:
\begin{equation}
  \log(\text{TotalWages}_{i}) =
    \beta_0 + \beta_1 \cdot \text{post2022}_{i}
    + \mathbf{\beta}_2 \cdot \text{Rank}_i
    + \mathbf{\beta}_3 \cdot \text{Campus}_i
    + \beta_4 \cdot \text{Year}_i + \varepsilon_i,
  \label{eq:baseline}
\end{equation}
where $\text{post2022}$ is an indicator equal to 1 for observations from
2022 onward and 0 otherwise. $\text{Rank}$ and $\text{Campus}$ enter as
sets of dummy variables (fixed effects), with Assistant Professor and UC
Berkeley as the respective reference categories. $\text{Year}$ is a
continuous linear time trend that controls for secular growth in nominal
wages. The coefficient $\beta_1$ captures whether log wages deviated
from the underlying trend after 2022, conditional on rank and campus.
Using a log transformation on the outcome allows coefficients to be
interpreted as approximate percentage changes.

Because the wage distribution is right-skewed and exhibits
heteroskedasticity---variance among Full Professors is substantially
wider than among Assistant Professors---we use heteroskedasticity-robust
(HC1) standard errors throughout. We verify that multicollinearity
between the continuous $\text{Year}$ trend and the binary
$\text{post2022}$ step function is not severe by computing Variance
Inflation Factors (VIF).

As a robustness check, we re-estimate Equation~\eqref{eq:baseline} with
$\log(\text{RegularPay})$ as the dependent variable. Regular pay reflects
only base compensation, excluding bonuses and other discretionary pay.
Consistency across both outcomes confirms that findings are not an
artifact of variation in bonus structures.

We also run the same specification separately within each rank---the
\textit{within-rank regressions}---which hold rank composition constant
by construction and test whether the post-2022 level shift persists
inside individual faculty categories.

\subsection{Levene's Test for Variance Equality}

We test whether the variance of salaries changed significantly between
the pre-2022 period (2012--2021) and the post-2022 period (2022--2024)
within each rank using Levene's test for equality of variances.
Levene's test is robust to non-normal distributions, making it
appropriate for wage data. The variance ratio (post-2022 variance divided
by pre-2022 variance) indicates the direction of change: a ratio greater
than 1 indicates that absolute variance \emph{increased} after 2022,
while a ratio less than 1 indicates it decreased.

\subsection{Difference-in-Differences Interaction Model}

To test whether wage gaps \emph{between} ranks changed after 2022---the
strongest form of the compression hypothesis---we estimate:
\begin{equation}
  \log(\text{TotalWages}_{i}) =
    \gamma_0 + \gamma_1 \cdot \text{post2022}_{i}
    \times \text{Rank}_i
    + \mathbf{\gamma}_2 \cdot \text{Campus}_i
    + \gamma_3 \cdot \text{Year}_i + \varepsilon_i.
  \label{eq:did}
\end{equation}
The interaction terms $\text{post2022} \times \text{Rank}$ capture
whether the post-2022 wage shift was differentially larger or smaller
for a given rank relative to the reference category (Assistant
Professor). A negative interaction for Full Professors would indicate
that Full Professors experienced a larger proportional decline than
Assistant Professors after 2022, consistent with compression from the
top of the distribution.

We note an important limitation of this design: our data are a repeated
cross-section with no individual identifiers. The DiD model therefore
captures changes in the \emph{distribution} of wages across groups, not
individual wage trajectories. We cannot rule out unobserved compositional
shifts within ranks---for example, if the mix of specializations or
seniority levels within a rank changed after 2022.

% ============================================================
\section{Results}
\label{sec:results}

\subsection{Salary Trends}
\label{subsec:trends}

Figure~\ref{fig:salary_trends} displays the evolution of mean total
wages, year-over-year growth, quartile salary trends, and total professor
headcount from 2012 to 2024. Mean nominal wages grew steadily and
monotonically from 2012 through 2023, with the 2024 observation
representing the first nominal decline in the sample period ($-6.4$\,\%).
The IQR widened over time, reflecting growing absolute dispersion, but
the rate of widening decelerated after 2022 for most ranks. The
professor headcount grew from roughly 16,000 in 2012 to just over 20,000
in 2023, with a decline to approximately 18,900 in 2024.

Figure~\ref{fig:headcount_by_rank} shows headcount by rank over time.
The Full Professor category is the largest throughout the sample. The
Clinical Professor category grew substantially, from roughly 1,400
observations in 2012 to over 2,500 by 2024, reflecting expanded
health-system activity across UC campuses.

\subsection{Within-Rank Coefficient of Variation}
\label{subsec:cv}

Figure~\ref{fig:cv_table} and Figure~\ref{fig:cv_by_rank} report the CV for
each rank and year. The average pre-2022 CV (2012--2021 mean) and
post-2022 CV (2022--2024 mean) for each rank are as follows:

\begin{table}[H]
\centering
\caption{CV Summary: Pre-2022 vs.\ Post-2022 Average}
\label{tab:cv_summary}
\begin{tabular}{lrrr}
\toprule
Rank & Avg CV Pre-2022 (\%) & Avg CV Post-2022 (\%) & Change (pp) \\
\midrule
Assistant  & 73.13 & 74.55 & $+1.42$ \\
Associate  & 76.52 & 71.90 & $-4.61$ \\
Full       & 59.22 & 56.30 & $-2.92$ \\
Clinical   & 64.81 & 64.19 & $-0.63$ \\
Research   & 68.53 & 65.64 & $-2.89$ \\
\bottomrule
\end{tabular}
\begin{flushleft}
\small Pre-2022 = average CV across 2012--2021; Post-2022 = average CV
across 2022--2024. CV computed as $100 \times \sigma / \mu$ within each
rank-year cell. Source: notebook cell 36 (deduplicated regression panel, N\,=\,229,386).
\end{flushleft}
\end{table}

Four of the five ranks show a decline in relative wage dispersion after
2022. The largest drop occurs among Associate Professors
($-4.61$ percentage points), followed by Full Professors ($-2.92$\,pp)
and Research Professors ($-2.89$\,pp). Clinical Professors exhibit a
modest decline of $-0.63$\,pp. Assistant Professors are the notable
exception: their CV increased by $+1.42$\,pp, indicating that relative
wage dispersion among entry-level faculty \emph{widened} after 2022.

The CV trends for Full and Associate Professors show a broadly
monotonic decline from 2012 onward that predates the 2022 threshold,
suggesting that within-rank convergence was a gradual ongoing process
rather than a discrete policy response. The 2022 threshold does not
appear to coincide with a sudden break for most ranks. The Assistant
Professor CV, in contrast, declined through 2020 before reversing
upward in 2022--2024.

\subsection{Levene's Test for Variance Equality}
\label{subsec:levene}

Figure~\ref{fig:levene} reports Levene's test
results. Across all five ranks, the null hypothesis of equal variances
is rejected at $p < 0.001$, indicating that the variance of salaries
changed significantly between the pre- and post-2022 periods. However,
the variance \emph{ratios} range from 1.32 (Full) to 1.59 (Assistant),
all greater than 1, indicating that absolute salary variance
\emph{increased} post-2022 across every rank.

This result appears to contradict the CV compression finding but is
fully reconcilable. Mean salaries rose substantially during the
post-2022 period, driven by inflation, secular wage growth, and the
consequences of the 2022 labor action. Because mean salaries grew faster
than the absolute spread of the distribution for most ranks, relative
dispersion (as measured by the CV) fell even as the dollar-denominated
spread widened. The appropriate measure for assessing the compression
hypothesis is therefore the CV, not the raw variance: workers experience
inequality in proportional terms relative to their peers' wages, not in
absolute dollar differences.

\subsection{Baseline OLS Results}
\label{subsec:baseline}

Figure~\ref{fig:table_tw} presents the
baseline OLS results for $\log(\text{TotalWages})$. The post-2022
coefficient is $-0.0634$ (SE $= 0.0063$, $p < 0.001$), indicating that
log total wages were approximately 6.3\,\% below the level predicted by
the secular trend after 2022, conditional on rank and campus. All rank
fixed effects are positive and strongly significant relative to the
Assistant Professor baseline, confirming a clear compensation hierarchy:
Clinical $>$ Full $>$ Research $>$ Associate $>$ Assistant. The linear
time trend coefficient of $0.0475$ implies annual nominal wage growth of
roughly 4.9\,\% per year along the pre-trend.

Figure~\ref{fig:table_rp} presents the
robustness check using $\log(\text{RegularPay})$ as the dependent
variable. The post-2022 coefficient is $-0.0410$ (SE $= 0.0043$,
$p < 0.001$). The smaller magnitude relative to the total-wages result
suggests that some of the post-2022 decline in total wages reflects
reduced variable pay (bonuses, overload pay), but the direction and
significance of the finding are unchanged. Both models use the
deduplicated panel of 229,386 and 225,412 observations respectively.

\subsection{Multicollinearity Check}
\label{subsec:vif}

Figure~\ref{fig:vif} reports the Variance Inflation Factors (VIF) for
the baseline regression predictors. The VIF for $\text{Year}$ is 2.197
and for $\text{post2022}$ is 2.193, both well below the commonly used
concern threshold of 5 and far below the severe threshold of 10.
The VIF for rank and campus are approximately 1.0. Multicollinearity
is not a concern in the baseline specification.

\subsection{Difference-in-Differences Interaction Results}
\label{subsec:did}

Figure~\ref{fig:table_did} presents the
interaction regression results. The baseline $\text{post2022}$ term
is $-0.0813$ (SE $= 0.0103$, $p < 0.001$), reflecting the post-period
wage decline for Assistant Professors (the reference group).

The interaction term for Full Professors is $-0.0016$ ($p = 0.888$),
which is economically negligible and statistically indistinguishable
from zero. This means Full Professors experienced essentially the same
post-2022 wage shift as Assistant Professors, providing no evidence of
between-rank compression from the top.

Associate Professors show a statistically significant positive
interaction ($\text{coef} = +0.0476$, $p = 0.0002$), indicating they
fared \emph{better} relative to Assistant Professors after 2022. This
is broadly consistent with the strong CV compression observed among
Associates: their wages held up better than those of either Assistant or
Full Professors in the post-period. Clinical Professors also show a
positive interaction ($\text{coef} = +0.0548$, $p = 0.002$), reflecting
continued compensation growth in health-system faculty markets. The
Research Professor interaction is positive but statistically
insignificant ($\text{coef} = +0.0301$, $p = 0.253$).

Taken together, the DiD evidence does not support the hypothesis that
the UC systematically narrowed the wage gap between senior and junior
faculty after 2022 through targeted compression at the top.

\subsection{Within-Rank Regressions}
\label{subsec:withinrank}

Figures~\ref{fig:wr_assistant}--\ref{fig:wr_research} report the
within-rank baseline regressions, which hold rank composition constant
and directly test whether post-2022 wages deviated from the within-rank
secular trend. Results are summarized in Table~\ref{tab:within_rank}.

\begin{table}[H]
\centering
\caption{Within-Rank Post-2022 Coefficient Summary}
\label{tab:within_rank}
\begin{tabular}{lrrrrl}
\toprule
Rank & Coef. & Std.\ Err. & $t$ & $p$-value & Approx.\ Effect \\
\midrule
Assistant  & $-0.0954$ & 0.0132 & $-7.23$ & $<0.001$ & $-9.1\%$ \\
Associate  & $-0.0452$ & 0.0130 & $-3.47$ & $0.001$  & $-4.4\%$ \\
Full       & $-0.0728$ & 0.0086 & $-8.49$ & $<0.001$ & $-7.0\%$ \\
Clinical   & $-0.0272$ & 0.0235 & $-1.16$ & $0.247$  & $-2.7\%$ (n.s.) \\
Research   & $+0.0092$ & 0.0375 & $+0.25$ & $0.805$  & $+0.9\%$ (n.s.) \\
\bottomrule
\end{tabular}
\begin{flushleft}
\small Specification: $\log(\text{TotalWages}) \sim
\text{post2022} + C(\text{EmployerName}) + \text{Year}$, HC1 robust
SE. Approximate percentage effect computed as $(\exp(\text{coef})-1)
\times 100$.
\end{flushleft}
\end{table}

The within-rank results reveal an important pattern: the post-2022
wage shortfall relative to trend is largest for Assistant Professors
($-9.1\%$, $p < 0.001$) and Full Professors ($-7.0\%$, $p < 0.001$),
while it is smaller and insignificant for Clinical and Research
professors. Associate Professors show a moderate, significant decline of
$-4.4\%$. The fact that Assistant Professors show the largest
within-rank decline in absolute magnitude, despite their CV rising,
reflects that the \emph{level} of their wages fell relative to trend
while the \emph{spread} of their distribution widened---consistent with
the external market setting floor wages that vary widely across
disciplines.

\subsection{Pre-versus Post-2022 Salary Distributions}
\label{subsec:distributions}

Figure~\ref{fig:pre_post_dist} displays overlapping salary histograms for
each rank, comparing the pre-2022 pooled distribution to the post-2022
pooled distribution. For Full and Associate Professors, the post-2022
distribution shifts rightward, reflecting higher mean wages, while the
relative spread of the distribution narrows modestly. For Assistant
Professors, the rightward shift is present but the tails of the
distribution widen proportionally, consistent with the CV increase for
that rank. Clinical Professors show a broad rightward shift with a
pronounced right tail.

% ============================================================
\section{Discussion}
\label{sec:discussion}

\subsection{Interpreting the Main Findings}

The results present a nuanced picture of how the UC faculty labor market
responded to the 2022 transparency and labor-action shock. Two findings
stand out.

First, there is clear evidence of \emph{within-rank} relative wage
compression among established faculty. The CVs for Full, Associate,
Research, and Clinical professors all declined after 2022, with Associate
professors exhibiting the strongest signal ($-4.61$\,pp). This pattern
is consistent with the theoretical prediction in \citet{cullen2022}: when
wage comparisons become more salient, employers face pressure to narrow
within-group pay gaps to limit dissatisfaction among employees who
directly compare their compensation with same-rank peers.

Second, there is no evidence of \emph{between-rank} compression. The DiD
interaction for Full versus Assistant Professors is economically small
and statistically indistinguishable from zero. This is theoretically
coherent: inequality aversion operates among direct peers, not across
structurally distinct positions. An Assistant Professor has no basis for
comparing their salary to a Full Professor's, given the evident
differences in seniority, tenure, and accumulated human capital. The
utility specification in Equation~\eqref{eq:utility} implies that
$\lambda$ is activated by comparisons within a relevant reference group,
not across all employees.

The exception among ranks---Assistant Professors---highlights an
important boundary condition on transparency-driven compression. Entry
salaries are set through a competitive national academic job market with
substantial cross-disciplinary variation, and the UC cannot compress
entry wages without risking recruitment failures in high-demand fields.
The $+1.42$\,pp increase in the Assistant Professor CV is consistent
with growing salary heterogeneity in external markets for new faculty.

\subsection{Reconciling CV and Levene Results}

Levene's tests show that absolute salary variance increased significantly
after 2022 across all ranks, which superficially contradicts the
CV-based compression finding. The reconciliation is straightforward.
Mean salaries rose substantially during the post-2022 period, driven by
inflation, secular wage growth, and the historic 2022 academic strikes.
Because the mean grew faster than the standard deviation for senior
ranks, the CV declined even as the dollar-denominated spread of the
distribution widened. This illustrates why the CV is the appropriate
dispersion metric for the compression hypothesis: workers experience
inequality in proportional terms relative to their peers' compensation
levels, not as an absolute dollar gap.

\subsection{Clinical Professors}

The positive and significant DiD interaction for Clinical Professors
($\text{coef} = +0.0548$, $p = 0.002$) indicates they fared better than
Assistant Professors in the post-2022 period. This likely reflects
continued growth in health-system compensation at UC's medical campuses,
driven by broader dynamics in the physician and academic medicine labor
market. Clinical faculty operate under different compensation structures
than academic-track faculty, limiting the degree to which internal UC
equity norms constrain their pay.

\subsection{The 2024 Decline}

The $-6.4\%$ year-over-year decline in mean total wages in 2024 is the
most pronounced single-year change in the sample. At the same time, total
observations fell from 20,176 in 2023 to 18,898 in 2024, a decline of
approximately 6.3\%. Because our data are a repeated cross-section, we
cannot determine the extent to which the mean wage decline reflects
actual pay cuts versus compositional turnover---specifically, whether
higher-paid senior faculty exited the UC system while lower-paid new
hires entered. The 2024 pattern warrants cautious interpretation.

\subsection{Limitations}

\paragraph{Repeated cross-section.} The fundamental limitation of this
study is the absence of individual identifiers. We cannot track any
professor across years, so we cannot measure individual wage growth,
distinguish pay cuts from turnover, or control for unobserved
compositional shifts within ranks. It is possible that apparent
within-rank compression among Full Professors reflects adverse selection:
if the UC compressed wages by capping high earners, and those
individuals subsequently departed for private universities, the observed
compression in the remaining sample would reflect attrition rather than
deliberate wage-setting policy.

\paragraph{Identification.} Our design compares the post-2022 period to
the pre-2022 secular trend. We cannot claim that any observed change
\emph{caused} by the transparency mandate or the strikes, as multiple
forces changed simultaneously in 2022--2023. The post-2022 indicator
absorbs any event occurring at that threshold. We therefore frame our
findings as descriptive associations rather than causal estimates.

\paragraph{Nominal wages only.} All analysis uses nominal wages. Without
verified inflation data, we do not convert to real wages. Over a 13-year
period spanning significant inflation, the nominal trend likely
overstates real wage growth. However, our main finding concerns
\emph{relative} within-rank dispersion (the CV), which is unaffected by
common inflation because it divides the standard deviation by the mean.

% ============================================================
\section{Conclusion}
\label{sec:conclusion}

This paper examined within-rank salary dispersion among University of
California faculty using 13 years of administrative compensation data
spanning 2012--2024. Our analysis focused on whether the period
surrounding the 2022 California Pay Transparency Act and the historic UC
academic labor strikes corresponds with measurable changes in the
distribution of wages within faculty ranks.

Our main findings are as follows. Mean nominal faculty wages grew by
approximately 52\,\% over the sample period before declining 6.4\,\% in
2024. The coefficient of variation declined after 2022 for Full
($-2.9$\,pp), Associate ($-4.6$\,pp), Research ($-2.9$\,pp), and
Clinical ($-0.6$\,pp) professors, indicating a narrowing of relative pay
disparities within these ranks. Assistant Professor CV increased by
$+1.4$\,pp, consistent with external labor-market pressures that prevent
internal compression at the hiring margin. The baseline OLS regression
finds that log wages were approximately 6.3\,\% below the secular trend
after 2022, conditional on rank and campus. The DiD interaction model
finds no evidence that wage gaps between Full and Assistant Professors
narrowed, but Associate and Clinical professors fared relatively better
in the post-period. Levene's tests confirm that absolute variance
increased across all ranks, but this is consistent with CV compression
because mean salaries grew faster than the spread for senior ranks.

These findings suggest that wage transparency and labor actions can
induce relative salary compression \emph{within} established peer groups
but have a more limited impact \emph{across} structurally distinct
positions. Internal equity pressures appear to operate primarily among
direct peers who observe and compare their wages, while the
compensation of entry-level faculty remains tethered to external market
conditions. Future research with individual-level panel data would be
needed to separate deliberate wage-setting compression from compositional
turnover effects.

% ============================================================
\newpage
\bibliographystyle{apalike}
\begin{thebibliography}{99}

\bibitem[Akerlof(1970)]{akerlof1970}
Akerlof, G. A. (1970).
\newblock The market for ``lemons'': Quality uncertainty and the market
mechanism.
\newblock \textit{The Quarterly Journal of Economics}, 84(3), 488--500.

\bibitem[Bandiera et al.(2005)]{bandiera2005}
Bandiera, O., Barankay, I., \& Rasul, I. (2005).
\newblock Social preferences and the response to incentives: Evidence from
personnel data.
\newblock \textit{The Quarterly Journal of Economics}, 120(3), 917--962.

\bibitem[Card et al.(2012)]{card2012}
Card, D., Mas, A., Moretti, E., \& Saez, E. (2012).
\newblock Inequality at work: The effect of peer salaries on job satisfaction.
\newblock \textit{American Economic Review}, 102(6), 2981--3003.

\bibitem[Cullen(2024)]{cullen2024}
Cullen, Z. B. (2024).
\newblock Is pay transparency good?
\newblock \textit{Journal of Economic Perspectives}, 38(1), 153--180.

\bibitem[Cullen \& Pakzad-Hurson(2022)]{cullen2022}
Cullen, Z. B., \& Pakzad-Hurson, B. (2022).
\newblock Equilibrium effects of pay transparency.
\newblock \textit{Econometrica}, 91(3), 765--802.
% TODO: Verify final publication details; the NBER Working Paper version
% is No. 28903; the published version appeared in Econometrica (2023).

\bibitem[University of California(2025)]{UC2025}
University of California. (2025).
\newblock \textit{UC Employee Headcount}.
\newblock University of California Information Center.
\newblock Retrieved from \url{https://www.universityofcalifornia.edu/infocenter}

\end{thebibliography}

% ============================================================
\newpage
\appendix

\section*{Appendix A: Mathematical Model of Inequality Aversion}
\addcontentsline{toc}{section}{Appendix A: Mathematical Model}
\label{app:model}

To formalize the theoretical mechanism, we employ the inequality aversion
framework from \citet{bandiera2005}. The agent's utility function is:
\begin{equation}
  U_i = W_i - C(e_i) - \lambda\max(0,\; W_{\text{peer}} - W_i),
  \label{eq:util_app}
\end{equation}
where $W_i$ is the worker's wage, $C(e_i)$ is the convex cost of effort
$e_i$, and the final term captures the disutility from observing a
higher-paid peer, scaled by the salience parameter $\lambda \geq 0$.

Prior to 2022, UC salary data were publicly posted but received limited
attention; we model this as $\lambda \approx 0$, so pay comparisons
generated negligible disutility. The 2022 events---the Pay Transparency
Act, the strikes, and intensified media coverage of faculty pay---shifted
$\lambda > 0$ for a meaningful share of the workforce. Once $\lambda > 0$,
the incentive-compatibility constraint for effort is tightened: the
principal (the UC system) must either raise $W_i$ toward $W_{\text{peer}}$
or reduce $W_{\text{peer}}$ toward $W_i$ to maintain effort levels. Both
channels narrow the distribution. The result is that the optimal contract
involves lower within-group wage variance, i.e., wage compression.

The prediction is limited to comparisons within a relevant reference
group. An Assistant Professor does not compare their wage to a Full
Professor's because the rank difference is institutionally recognized and
perceived as structurally justified. The model therefore predicts
compression within ranks but not across ranks---consistent with our
empirical findings.

\vspace{1em}

\section*{Appendix B: Figures}
\addcontentsline{toc}{section}{Appendix B: Figures}
\label{app:figures}

\begin{figure}[H]
  \centering
  \includegraphics[width=\textwidth]{figures/fig_salary_trends.png}
  \caption{UC Professor Salary Trends, 2012--2024 (Nominal). Top left:
  mean total wages with $\pm$1 standard deviation band. Top right:
  year-over-year percentage change in mean salary; the 2024 bar
  ($-6.4$\,\%) is the only negative value in the sample. Bottom left:
  quartile salary trends (Q25, median, Q75). Bottom right: total
  professor observation count by year. Vertical dashed line marks the
  2022 threshold.}
  \label{fig:salary_trends}
\end{figure}

\begin{figure}[H]
  \centering
  \includegraphics[width=\textwidth]{figures/fig_cv_by_rank.png}
  \caption{Salary Dispersion by Professor Rank, 2012--2024. Top left:
  coefficient of variation (CV\%) by rank; lower values indicate less
  relative dispersion. Top right: standard deviation of salary by rank.
  Bottom left: interquartile range by rank. Bottom right: bar chart
  summarizing the CV change in percentage points (post-2022 minus
  pre-2022 average). Red bar indicates CV expansion (Assistant); green
  bars indicate CV compression. Vertical dashed line marks the 2022
  threshold.}
  \label{fig:cv_by_rank}
\end{figure}

\begin{figure}[H]
  \centering
  \includegraphics[width=\textwidth]{figures/fig_pre_post_distributions.png}
  \caption{Salary Distribution by Rank: Pre-2022 vs.\ Post-2022 (Total
  Wages, Nominal). Each panel overlays the pooled pre-2022 (blue) and
  post-2022 (red) salary histograms for one rank. Vertical dashed lines
  mark the mean for each period. For Full and Associate Professors, the
  post-2022 distribution shifts rightward while the relative spread
  narrows modestly. For Assistant Professors, the distribution widens
  proportionally with the rightward shift.}
  \label{fig:pre_post_dist}
\end{figure}

\begin{figure}[H]
  \centering
  \includegraphics[width=0.85\textwidth]{figures/fig_headcount_by_rank.png}
  \caption{Professor Headcount by Rank, 2012--2024. Full Professors
  (green) are consistently the largest group. Clinical Professors (red)
  grew substantially from roughly 1,400 in 2012 to over 2,500 by 2024.
  Assistant Professors (blue) grew through 2019--2020 before declining
  in 2023--2024. Vertical dashed line marks the 2022 threshold.}
  \label{fig:headcount_by_rank}
\end{figure}

\vspace{1em}

\section*{Appendix C: Tables}
\addcontentsline{toc}{section}{Appendix C: Tables}
\label{app:tables}

\begin{figure}[H]
  \centering
  \includegraphics[width=0.9\textwidth]{figures/table_cv_by_rank_year.png}
  \caption{Coefficient of Variation (CV\%) by Rank and Year, 2012--2024.
  Computed as $100 \times \sigma / \mu$ within each rank-year cell on
  the deduplicated panel. Sample restricted to professors with
  TotalWages $> 0$.}
  \label{fig:cv_table}
\end{figure}

\begin{figure}[H]
  \centering
  \includegraphics[width=0.9\textwidth]{figures/table_baseline_total_wages.png}
  \caption{Table 1: Baseline OLS regression of log(Total Wages) on
  post2022, rank fixed effects, campus fixed effects, and a linear year
  trend. $N = 229{,}386$. HC1 heteroskedasticity-robust standard errors.
  Reference category: Assistant Professor.}
  \label{fig:table_tw}
\end{figure}

\begin{figure}[H]
  \centering
  \includegraphics[width=0.9\textwidth]{figures/table_baseline_regular_pay.png}
  \caption{Table 2: Robustness check---baseline OLS regression of
  log(Regular Pay) on post2022, rank fixed effects, campus fixed effects,
  and a linear year trend. $N = 225{,}412$. HC1 robust standard errors.
  Reference category: Assistant Professor.}
  \label{fig:table_rp}
\end{figure}

\begin{figure}[H]
  \centering
  \includegraphics[width=0.9\textwidth]{figures/table_did_compression.png}
  \caption{Table 3: Difference-in-differences compression test. Dependent
  variable: log(Total Wages). Specification includes post2022, the full
  interaction post2022 $\times$ Rank, campus fixed effects, and a linear
  year trend. $N = 229{,}386$. HC1 robust standard errors. Reference
  category: Assistant Professor.}
  \label{fig:table_did}
\end{figure}

\begin{figure}[H]
  \centering
  \includegraphics[width=0.9\textwidth]{figures/table_levene_test.png}
  \caption{Levene's Test for Variance Equality---Pre-2022 vs.\
  Post-2022. Variance Ratio $> 1$ indicates that absolute variance
  \emph{increased} post-2022. All five ranks show statistically
  significant variance changes ($p < 0.001$). This result is consistent
  with CV compression because mean salaries rose faster than the absolute
  spread for senior ranks.}
  \label{fig:levene}
\end{figure}

\begin{figure}[H]
  \centering
  \includegraphics[width=0.9\textwidth]{figures/table_vif.png}
  \caption{Variance Inflation Factors (VIF) for the baseline regression
  predictors. All VIFs are well below the concern threshold of 5,
  confirming that multicollinearity between the continuous Year trend and
  the binary post2022 indicator is not severe.}
  \label{fig:vif}
\end{figure}

\begin{figure}[H]
  \centering
  \includegraphics[width=0.9\textwidth]{figures/table_within_rank_assistant.png}
  \caption{Within-Rank Regression: Assistant Professors. Specification:
  $\log(\text{TotalWages}) \sim \text{post2022} + C(\text{EmployerName})
  + \text{Year}$. $N = 63{,}756$. HC1 robust standard errors.}
  \label{fig:wr_assistant}
\end{figure}

\begin{figure}[H]
  \centering
  \includegraphics[width=0.9\textwidth]{figures/table_within_rank_associate.png}
  \caption{Within-Rank Regression: Associate Professors.
  $N = 50{,}988$. Specification and standard errors as in
  Figure~\ref{fig:wr_assistant}.}
  \label{fig:wr_associate}
\end{figure}

\begin{figure}[H]
  \centering
  \includegraphics[width=0.9\textwidth]{figures/table_within_rank_full.png}
  \caption{Within-Rank Regression: Full Professors.
  $N = 83{,}141$. Specification and standard errors as in
  Figure~\ref{fig:wr_assistant}.}
  \label{fig:wr_full}
\end{figure}

\begin{figure}[H]
  \centering
  \includegraphics[width=0.9\textwidth]{figures/table_within_rank_clinical.png}
  \caption{Within-Rank Regression: Clinical Professors.
  $N = 23{,}918$. The post2022 coefficient is negative but not
  statistically significant ($p = 0.247$), reflecting the distinct
  dynamics of health-system compensation.}
  \label{fig:wr_clinical}
\end{figure}

\begin{figure}[H]
  \centering
  \includegraphics[width=0.9\textwidth]{figures/table_within_rank_research.png}
  \caption{Within-Rank Regression: Research Professors.
  $N = 7{,}583$. The post2022 coefficient is positive and not
  statistically significant ($p = 0.805$), indicating no significant
  deviation from trend within this smaller rank category.}
  \label{fig:wr_research}
\end{figure}

\end{document}
